

















































































































































\section{Introduction}
We study a class of optimization problems ubiquitous in theoretical computer science, machine learning,  quantum information theory and statistics. The programs we consider are continuous optimization problems over matrix groups. More precisely, they can be posed as Euclidean norm minimization over the closure of a group orbit. The programs span two historically distinct contexts: In one context, the optimization problems are convex, and in the other they are not convex but rather \emph{geodesically convex} on a suitable manifold.


The \emph{commutative setting}, in which the underlying group is Abelian, captures matrix scaling, matrix balancing and array scaling, which arise in scientific computing and optimal transport \cite{cuturi2013sinkhorn, parlett1971balancing}. Such problems fall into the framework of unconstrained geometric programming. Though these problems are convex, there are at least two reasons to study them further. Firstly, they are of such practical importance that speed matters. Na\"ively applying powerful algorithms like ellipsoid and interior point methods can be impractically slow. Hence, it is important to understand when faster methods can succeed. Matrix scaling and balancing, in particular, have enjoyed some success stories~- there are fast algorithms to obtain high precision solutions \cite{cohen2017matrix,allen2017much}, and there are more general upper bounds \cite{burgisser2020interior}. Secondly, the algorithms developed for the commutative setting are candidates for generalization to our second setting, which takes place in the less well-understood arena of geodesically convex optimization.


The second context, which we call the \emph{noncommutative setting}, arises when the underlying group is non-Abelian. The noncommutative setting captures problems like operator and tensor scaling \cite{garg2016deterministic,burgisser2017alternating},  the quantum marginal problem \cite{burgisser2018efficient} and statistical estimators such as Tyler's M estimator \cite{franks2020rigorous} and maximum likelihood estimates for matrix and tensor normal models \cite{amendola2020invariant}. Deciding whether the value of the optimization problem is zero or not is equivalent to deciding a central polynomial identity testing (P.I.T.) problem in invariant theory known as the \emph{null cone} problem. It is hoped that efficient optimization algorithms will result in efficient algorithms for the null-cone problem.
One approach to complexity lower bounds, geometric complexity theory, suggests that these P.I.T. problems should be in $\P$ \cite{mulmuley2017geometric, garg2019search}, and the optimization approach has resulted in polynomial time algorithms in some cases \cite{garg2016deterministic, allen2018operator}. The optimization problems that arise in the noncommutative setting are not convex in the Euclidean sense, but rather \emph{geodesically convex}, a notion of convexity on a Riemannian manifold. Currently, the only implementable algorithms for geodesically convex optimization are analogues of gradient descent and trust region methods \cite{absil2008optimization, zhang2016first,allen2018operator}. There are, as of yet, no efficiently implementable geodesically convex counterparts to the interior point or cutting plane methods.


In both the commutative and noncommutative settings, algorithms are typically analysed using two quantities. One is \emph{diameter}, or how far approximate minimizers can be from the origin. The other is a geometric measure of well-conditionedness known as \emph{margin} (or \emph{gap} in the noncommutative case), which has several variants in the literature and appears in two primary ways. Firstly, the smaller the margin, the higher the degree of precision required to decide if the value of the optimization problem is zero or not \cite{gradflow, gurvits2004combinatorial}. Secondly, the larger the margin, the smaller the diameter \cite{singh2014entropy,straszak2017computing,gradflow,burgisser2020interior}. In this paper we show the following:


\begin{enumerate} 
\item[i)] In the commutative setting, and in particular for array scaling, approximate minimizers for the functions we study can have doubly exponential condition number. That is, the problems have exponential diameter. As a consequence, popular classes of algorithms such as gradient descent and trust region methods cannot produce high-precision solutions in polynomial time in general. This result applies in the noncommutative setting as well, which provides evidence that even cutting plane methods are unlikely to produce high-precision solutions in polynomial time. This shows it is necessary to develop powerful methods like the interior point method in the geodesically convex setting. 
\item[ii)] In the commutative and noncommutative settings, we study the margin and gap, respectively, which appear in running time bounds for all existing algorithms. We prove that these measures can be exponentially small in the input size for several problems including array scaling and tensor scaling. In the commutative case, this gives evidence that existing algorithms for array scaling do not run in near-linear time. In the noncommutative case, our results show that margin-based analyses like \cite{gradflow} cannot prove polynomial time guarantees for deciding the null cone problem for tensor scaling using trust region methods. 
\end{enumerate} 


We use the remainder of the introduction to describe both settings in more detail, state our main results precisely, and discuss previous work. For both the commutative and noncommutative settings, we proceed in the following order. We start with an introduction and motivation of the setting, continue with diameter bounds and afterwards treat bounds on the margin and gap, respectively. We end each setting with a short discussion of the main proof techniques.



\subsection{The commutative setting: matrix scaling and its relatives}


\paragraph{Matrix scaling and array scaling.}

Consider the matrix scaling problem: given a nonnegative matrix $A$, find nonnegative diagonal matrices $X, Y$ such that $XA Y$ is doubly stochastic (i.e. has row and column sums equal to one). The matrices, if they exist, can be found by the exceedingly simple and fast alternating minimization method known as \emph{Sinkhorn's algorithm}. It is frequently used in practice, e.g. for quickly approximating the solution to optimal transport problems \cite{cuturi2013sinkhorn}.

Like all other algorithms for matrix scaling,  Sinkhorn's algorithm is typically analyzed through optimization. One finds that $X$ and $Y$ are $e^{\diag(x)}, e^{\diag(y)}$, where $x,y \in \RR^n$ are solutions to the following optimization problem: 
\begin{align}\inf_{x, y \in \RR^n} \sum A_{ij}e^{x_i + y_j - \bar{x} - \bar{y}}\label{eq:matrix-capacity}\end{align}
for $\bar{z}:=\frac{1}{n} \sum z_i$ (c.f. \cite{kalantari1996complexity}). Moreover, the infimum is greater than zero if and only if $A$ is \emph{approximately scalable}, i.e. the row and column sums of $XAY$ can be made arbitrarily close to one for $X, Y$ nonnegative, diagonal.


More generally, given a finite set $\Omega \subseteq \RR^m$ and a nonnegative function $p:\Omega \to \RR_{\geq 0}$, define the \emph{capacity} \cite{gurvits2004combinatorial} as the value of the unconstrained geometric program
\begin{gather}
\capa(p):= \inf_{x \in \RR^m} f_p(x) := \inf_{x \in \RR^m} \sum_{\omega \in \Omega} p_\omega e^{\omega \cdot x}.\label{eq:abelian}
\end{gather}
The capacity is positive if and only if zero is in the \emph{Newton polytope} $\conv(\supp p)$. Matrix scaling arises when $m = 2n$ and $\Omega = \{(\eps_i, \eps_j): i,j \in [n]\}$ for $\eps_k:=e_k - \frac{1}{n} \id_n$, where $e_k \in \RR^n$ is the $k^{th}$ canonical unit vector and $\id_n \in \RR^n$ denotes the all-ones vector. In this case \cref{eq:abelian} reduces to precisely \cref{eq:matrix-capacity}, and $\|\nabla \log f_p (x)\|$ measures the deviation of $p$ from doubly stochastic.

Matrix balancing, in which we instead wish to find a scaling for which the $i^{th}$ row and column sum match, arises when $m = n$ and $\Omega = \{e_i - e_j: i\neq j \in [n]\}$. When $m = 3n$ and $\Omega = \{(\eps_i, \eps_j, \eps_k): i,j,k \in [n]\}$ we obtain the $3$-dimensional \emph{array scaling problem}. In analogy to matrix scaling, in array scaling one has an array $p$ of numbers in $(\RR^n_{\geq 0})^{\ot 3}$ and seeks positive vectors $X,Y,Z \in \RR^n_{\geq 0}$ so that the array $q$ with entries $q_{ijk} =p_{ijk}X_i Y_j Z_k$ is \emph{tristochastic}. That is, the sum over every \emph{slice} is equal to one, i.e.
    $\sum_{j,k} q_{i_0,j,k} = \sum_{i,k} q_{i,j_0,k} = \sum_{i,j} q_{i,j,k_0} = 1$
for all $i_0,j_0,k_0 \in [n]$. If it is possible to satisfy these equations to arbitrary precision we say $p$ is approximately scalable. As for matrix scaling, $p$ is approximately scalable if and only if $\capa(p) > 0$. In the same manner, we obtain $d$-dimensional array scaling for $m=dn$ and
	\begin{equation}\label{eq:defnOmega-n-d}
	\Omega = \Omega_{n,d} := \big\lbrace \eps_i \colon i \in [n] \big\rbrace^d \subseteq \big( \RR^n \big)^d.
	\end{equation}
We can think of subsets of $\Omega_{n,d}$ as $d$-uniform, $d$-partite hypergraphs. Up to an additive shift by $- \frac{1}{n} \id_{nd}$, the elements of $\Omega_{n,d}$ are indicator vectors of the edges in such hypergraphs. For $d = 2$, the matrix $p$ is scalable if and only if the bipartite graph corresponding to $\supp p$ contains a perfect matching, but this is not the case for $d \geq 3$ (indeed, $d$-partite hypergraph matching is $\NP$-hard).



\paragraph{Algorithms for array scaling.}

Array scaling serves the same role for speeding up multimarginal transport as matrix scaling for optimal transport, and yet again there is a simple and fast alternating minimization algorithm that produces $\eps$-tristochastic scalings in time $O(1/\eps^2)$ \cite{altschuler2020polynomial,lin2019complexity}.
Moreover, algorithms to approximate the capacity arise in varied settings including radial isotropic position \cite{hardt2013algorithms}, entropy maximization \cite{straszak2017computing}, and approximate counting \cite{anari2018log}.  
 

It is natural to ask if there are \emph{high-precision} algorithms for array scaling with $\log(1/\eps)$ dependence on the error and linear or mild dependence on the number of nonzero entries. For matrix scaling and matrix balancing, several works have shown that trust regions and interior point methods \emph{can} obtain such guarantees \cite{cohen2017matrix,allen2017much}.
Our work is concerned with whether the performance of such algorithms carries over to array scaling and the computation of the capacity in general.





































\setcounter{thm}{1}\subsubsection{Margins: the geometry of scaling problems}
Many computational aspects of the capacity rely on the convex geometry of the finite set $\Omega \subseteq \RR^m$. Consider the following quantity, which we call the \emph{margin} of $\Omega$. 
The margin is the minimum \emph{positive} distance from a convex hull of a subset of $\Omega$ to the origin. Formally,
\begin{dfn}[Margin]\label{dfn:margin} For a finite set $\Omega \subseteq \RR^m$, define the margin $\gamma(\Omega)$ by
	\[ \gamma(\Omega) := \min \left\lbrace \dist \big(0, \conv(S) \big) \mid S \subseteq \Omega, \; 0 \notin \conv(S) \right\rbrace.\]
\end{dfn}
We point out that for all considered capacity problems in this paper, the margin is actually the \emph{weight margin} (c.f. \cite{gradflow} and our Definition 4.3 of the original publication) of a certain group representation. For example, the margin for array scaling is the weight margin for tensor scaling. 
We now discuss how the margin enters in decision problems and diameter bounds.


\paragraph{Margin as a precision parameter for the decision problem.}


To illustrate how the margin enters the decision problem of whether $\capa(p) > 0$, consider matrix scaling.
To certify that the capacity of a matrix is nonzero, we compute $\eps$-doubly stochastic scalings for some $\eps$ smaller than the distance to doubly stochastic attained by any matrix that is \emph{not} approximately scalable. This turns out to be precisely $\gamma(\Omega_{n,2})$. More generally, it is a classical fact that for $p$ with support contained in $\Omega$, the gradient $\nabla \log f_p(x)$ can take any value in the Newton polytope of $p$. Thus, $\capa(p) > 0$ if and only if there is some $x$ with $\|\nabla \log f_p(x)\| \leq \gamma(\Omega)$.

For matrix scaling and matrix balancing, it is known that $\gamma(\Omega)$ is on the order of $n^{-3/2}$, despite the exponential number of subsets $S \subseteq \Omega$! This luck can be attributed to the extraordinary geometry of $\Omega$ in these cases, whose elements form the rows of a totally unimodular matrix (up to a shift).
On the other hand, for $d$-dimensional array scaling for $n = 2$, the margin $\gamma(\Omega_{2,d})$ is on the order of the margin of the $d$-dimensional hypercube $\{\pm 1\}^d$, which satisfies $\gamma \big( \{\pm 1\}^d \big)=d^{-\frac{d}{2}(1 + o(1))}$ by \cite{alon1997anti}.
However, between the extreme cases $\Omega_{n,2}$ (matrix scaling) and $\Omega_{2,d}$ (the hypercube), very little is known.


\paragraph{Margin and related quantities for diameter bounds.}
In addition to their role in the decision problem, margins and related quantities can be used to prove diameter bounds for \cref{eq:abelian}. The work \cite{gradflow} proves the diameter bound $\poly(\gamma(\Omega)^{-1}, \log(1/\eps))$. In \cite{straszak2017computing} it is shown that the diameter is polynomial in the logarithm of the minimum nonzero $p_\omega$ and a quantity called the \emph{unary facet complexity}. The latter is defined as the maximal length of an integer normal vector of a face of the Newton polytope $\conv(\supp p)$. In the case of $d$-dimensional arrays, one can use Cramer's rule to crudely bound the unary facet complexity by $(d + 1)^{dn}$. In the case when $0$ is in the relative interior of the Newton polytope, \cite{singh2014entropy} has shown that there is a minimizer with Euclidean norm $O(\log |\supp p|/\eta)$, where $\eta$ is the distance from $0$ to the boundary of the Newton polytope. The diameter bounds in \cite{singh2014entropy,straszak2017computing} were used to design ellipsoid methods that are tractable even for $|\supp p|$ very large, and in \cite{burgisser2020interior} they were used to bound the running time of interior point methods.

\paragraph{Main theorem.}
One is led to ask if the margin remains large for array scaling when $d \geq 3$. We show that this is not the case. In fact, the margin becomes exponentially small in $nd$ for $d \geq 3$. What follows is stated in more detail later in \cref{thm:MarginTensor}.

\begin{thm}\label{thm:tensor-margin} Let $d \geq 3$ and $n \geq 2$. Let $\Omega_{n,d} = \{\eps_{i}: i \in [n]\}^d \subseteq (\RR^n)^d$, where $\eps_j:=e_j - \frac{1}{n} \id_n$. There exists a constant $C>0$, independent of $n$ and $d$, such that
$\gamma(\Omega_{n,d}) \leq 2^{-Cnd}.$ 
\end{thm}

That is, there are $d$-dimensional arrays $p \in (\RR^n_{\geq 0})^{\ot d}$ such that the $d$-tuple of marginals of $p$ is at distance at most $2^{ - Cnd}$ from $\frac{1}{n}(\id_n, \dots, \id_n)$, yet the support of $p$ does not admit an array with uniform marginals, i.e. $\capa(p) = 0$. We note that the support of the array $p$ we construct has $O(nd)$ elements.


























































































































\section{The geometry of commutative scaling problems}\label{sec:MarginComm}

The purpose of this section is to show the following theorem on the margin of $d$-dimensional array scaling. Recall that the latter arises for $\Omega_{n,d} := \{\eps_i: i \in [n]\}^d \subseteq (\RR^n)^d$.

\begin{thm}[Margin for array scaling]\label{thm:MarginTensor}
The margin of $\Omega_{n,d} \subseteq (\RR^n)^d$ is bounded as follows.
	\begin{itemize}
	\item[(a)] If $n=2$ and $d \geq 3$, then
	$
	\gamma \left(\Omega_{2,d} \right)  \leq  2^{-\frac{d}{2} + 1}.
	$
	\item[(b)] If $n \geq 3$ and $d = 3$, then $\gamma(\Omega_{n,3}) \leq 2^{-n+1}$.
	\item[(c)] If $n \geq 3$ and $d = 6 r - 3$ for some integer $r \geq 2$, then
		\begin{align*}
		\gamma(\Omega_{n,d})  \leq \frac{\sqrt{6}}{(n-1)\sqrt{r}} \; 2^{-r(n-1) + 1} 
		\leq 2^{- r (n-1) + 1} = 2^{- \frac{(d+3)(n-1)}{6} + 1}.
		\end{align*}
		\end{itemize}
\end{thm}
By ``padding'' the tensors appropriately, one sees that a bound for $\gamma(\Omega_{n,d})$ also applies to $\gamma(\Omega_{n,d+1})$ (see  \cref{prp:dTensorsPadding}). Combining this result with \cref{thm:MarginTensor} above implies \cref{thm:tensor-margin} from the introduction. The next three subsections each prove one of the parts of \cref{thm:MarginTensor}; the construction for part~(a) with $n = 2$ is slightly different and the construction for part (c), $d > 3$ builds on the one for part~(b), $d = 3$.

To prove the results, we will frequently use the following simple lemma. Recall that an \emph{affine linear combination} of $v_1, \dots, v_k \in \RR^m$ is $\lambda_1 v_1 + \dots + \lambda_k v_k$ for $\lambda_i \geq 0, \sum_{i = 1}^k \lambda_i = 1$. The affine hull $\aff(S)$ of a set $S\subset \RR^m$ is the set of all affine linear combinations of finite subsets of $S$, or equivalently the affine space (i.e. translate of a subspace) of lowest dimension containing $S$.

\begin{lem}\label{lem:convCombEps-i}
In $\RR^n$ we have
\begin{equation}\label{eq:SumEps-i}
\sum_{i=1}^n \frac{1}{n} \: \varepsilon_i = 0_n 
\end{equation}
and this is the only affine linear combination of $\eps_1,\ldots,\eps_n$ giving zero.
\end{lem}

\begin{proof}
One calculates directly that $\sum_i \frac{1}{n} \, \eps_i = 0_n$. To show uniqueness of this affine combination, we note that the vectors $e_2, \ldots, e_n, \id_n$ are linearly independent. Thus, $\eps_2, \ldots, \eps_n$ are linearly independent. On the other hand, $\eps_1,\ldots, \eps_n$ are linearly dependent. Therefore
	$
	\left\lbrace (\lambda_1,\ldots,\lambda_n) \in \RR^n \mid \sum_i \lambda_i \: \eps_i = 0_n \right\rbrace
	$
is a one-dimensional subspace of $\RR^n$, which yields the uniqueness of the affine linear combination.
\end{proof}



\subsection{Local dimension two: the hypercube}\label{sec:qubits}

In this subsection we prove part (a) of \cref{thm:MarginTensor} by showing that the margin of $\Omega_{2,d}$ is exponentially small in $d$. This follows from \cite{alon1997anti}, but we present a new construction which has the additional property of \emph{freeness}, which we discuss later in Section 4 of the original publication. Recall that 
	\begin{align*}
	\Omega_{2,d} =  \big\lbrace (\eps_{i_1}, \ldots, \eps_{i_d}) \mid i_1, \ldots, i_d \in [2] \big\rbrace \subseteq \big( \RR^2 \big)^d.
	\end{align*}
In the following we construct a subset of $\Omega_{2,d}$, which witnesses the exponentially small margin. For this, we construct a matrix with entries in $[2]$, and each row of the matrix will correspond to an element of $\Omega_{2,d}$. For example, the row $(1,2,2)$ would correspond to $(\eps_1, \eps_2, \eps_2) \in \Omega_{2,3}$. To do so, we begin with the matrices 
	\begin{align*}
	A_2 : = \begin{pmatrix} 1 & 1 \\ 2 & 1 \end{pmatrix} , \;
	B_1 := \begin{pmatrix} 1 & 1 \\ 2 & 2 \end{pmatrix} , \;
	B_2 := \begin{pmatrix} 1 & 2 \\ 2 & 2 \end{pmatrix} , \;
	B_3 := \begin{pmatrix} 2 & 1 \\ 1 & 1 \end{pmatrix} ,
	\end{align*}
and define recursively
	\begin{equation}\label{eq:defA2r}
	A_{2r+2} := \begin{pmatrix}
	  &  &  & B_1 \\ 
	  &  A_{2r} &  & \vdots \\ 
	  &  &  & B_1 \\ 
	B_2 & \cdots & B_2 & B_3
	\end{pmatrix} = 
	\begin{pmatrix}
	 A_2 & B_1 & \cdots & B_1 \\ 
	 B_2  & B_3 & \ddots & \vdots \\ 
	 \vdots & \ddots & \ddots& B_1 \\ 
	 B_2 & \cdots & B_2 & B_3
	\end{pmatrix}
	\end{equation}
for $r \geq 1$. \cref{fig:qubit-matrices} is supplied as a visualization aid.

\begin{figure}
$$ A_4 = \left(\begin{array}{c|c}
\cellcolor{black!30}\begin{array}{cc} * & * \\  & * \end{array}&\cellcolor{green!10} \begin{array}{cc} * & * \\  &  \end{array}\\
\hline
\cellcolor{red!10}\begin{array}{cc} * & \;\text{ }  \\  &  \end{array}& \cellcolor{blue!10}\begin{array}{cc}  & * \\ * & * \end{array}
 \end{array}\right),
 \quad 
 A_6 = \left(\begin{array}{c|c|c} 
 \cellcolor{black!30} \begin{array}{cc} * & * \\ & * \end{array}& \cellcolor{green!10}\begin{array}{cc} * & * \\  &  \end{array} & \cellcolor{green!10} \begin{array}{cc} * & * \\  &  \end{array} \\
 \hline
\cellcolor{red!10}\begin{array}{cc} * & \;\text{ }  \\  &  \end{array} & \cellcolor{blue!10}  \begin{array}{cc}  & * \\ * & * \end{array} &  \cellcolor{green!10}\begin{array}{cc} * & * \\  &  \end{array}  \\
\hline
\cellcolor{red!10}\begin{array}{cc} * & \;\text{ } \\  &  \end{array}& \cellcolor{red!10}\begin{array}{cc} * & \;\text{ } \\  &  \end{array} &  \cellcolor{blue!10}\begin{array}{cc}  & * \\ * & * \end{array} \\
\end{array}\right)
  $$
\caption{The positions of the ones in $A_4$ and $A_6$ are marked by $*$ in the following figure and the cells are colored according to whether they belong to $A_2, B_1, B_2$ or $B_3$.}\label{fig:qubit-matrices}
\end{figure}


We remark that the entry of $A_{2r}$ at position $(i,j)$ is independent of $r$ and denote it by $a(i,j)$. We set for $r \geq 1$
	\begin{align*}
	\Gamma_{2,2r} &:= \left\lbrace \left( \eps_{a(i,1)}, \eps_{a(i,2)}, \ldots, \eps_{a(i,2r)} \right) \mid i \in [2r] \right\rbrace \subseteq \Omega(\pi_{2,2r}) \subseteq \big( \RR^2 \big)^{2r},\\
	\Gamma_{2,2r+1} &:= \left\lbrace \left( \eps_{a(i,1)}, \eps_{a(i,2)}, \ldots, \eps_{a(i,2r)}, \eps_{\chi(i)} \right) \mid i \in [2r] \right\rbrace \subseteq \Omega(\pi_{2,2r+1}) \subseteq \big( \RR^2 \big)^{2r+1},
	\end{align*}
where $\chi \colon \NN \to \{1,2\}, \; i \mapsto i \mod 2$. That is, $\Gamma_{2,2r}$ is the subset of $\Omega_{2,2r}$ induced by the rows of $A_{2r}$ and $\Gamma_{2,2r+1}$ is obtained by alternatingly appending $\eps_1$ or $\eps_2$ to the $2r$-many elements of $\Gamma_{2,2r}$.
	
\begin{lem}\label{lem:affHullQubits}
For $r \geq 1$ it holds that $0 \notin \aff(\Gamma_{2,2r})$ and $0 \notin \aff(\Gamma_{2,2r+1})$.
\end{lem}

\begin{proof}
By construction, $0 \in \aff(\Gamma_{2,2r+1})$ implies $0 \in \aff(\Gamma_{2,2r})$, so it suffices to prove $0 \notin \aff(\Gamma_{2,2r})$. We proceed by induction on $r \geq 1$. For $r=1$, it is clear that $0 \notin \aff(\Gamma_{2,2}) \subseteq \RR^2 \times \lbrace \eps_1 \rbrace$.
Now assume that $0 \notin \aff(\Gamma_{2,2r})$. For the sake of contradiction, let
	\begin{equation}\label{eq:affCombGamma}
	\sum_{i=1}^{2r+2} \lambda_i \left( \eps_{a(i,1)}, \eps_{a(i,2)}, \ldots, \eps_{a(i,2r+2)} \right) = 0 \in \left( \RR^2 \right)^{2r+2}
	\end{equation}
be an affine linear combination of $\Gamma_{2,2r+2}$. Then equation~\eqref{eq:affCombGamma} gives in each of the $(2r+2)$-many $\RR^2$-components the affine linear combination $2^{-1} (\eps_1 + \eps_2) = 0$, by Lemma~\ref{lem:convCombEps-i}. Considering the scalar factor of $\eps_1$ in the first, the penultimate and the last $\RR^{2}$-component respectively, we conclude
	\begin{align*}
	\underbrace{ \sum_{j=1}^{r+1} \lambda_{2j-1} }_{\text{first}} = \frac{1}{2}
	= \underbrace{ \lambda_{2r+2} + \sum_{j=1}^{r} \lambda_{2j-1} }_{\text{penultimate}} 
	= \frac{1}{2} = \underbrace { \lambda_{2r+2} + \sum_{j=1}^{r+1} \lambda_{2j-1} }_{\text{last}}
	\end{align*}
by construction of $A_{2r+2}$. Hence, $\lambda_{2r+2} = 0$ using the first and last component. Furthermore, the first and penultimate column give $\lambda_{2r+1} = \lambda_{2r+2} = 0$. Therefore, the first $2r$-many components in \cref{eq:affCombGamma} show $0 \in \aff(\Gamma_{2,2r})$, which contradicts our induction hypothesis.
\end{proof}

\begin{lem}\label{lem:convCombQubits}
For $r \geq 1$ it holds that $\dist(0, \conv(\Gamma_{2,2r}) ) \leq 2^{-r+ \frac{1}{2}}$ and $\dist(0, \conv(\Gamma_{2,2r+1}) ) \leq 2^{-r+ \frac{1}{2}}$.
\end{lem}

\begin{proof}
We first prove the inequality for $\conv(\Gamma_{2,2r})$. For $i \in [2r]$ let $\omega_i := \big( \eps_{a(i,1)}, \ldots, \eps_{a(i,2r)} \big) \in \left( \RR^2 \right)^{2r}$ be the weight in $\Gamma_{2,2r}$ that corresponds to the $i^{th}$ row of $A_{2r}$. Consider the convex combination
	\begin{equation}\label{eq:convCombGamma}
	(x_1, \ldots, x_{2r}) := 2^{-r} ( \omega_{2r-1} + \omega_{2r} ) + \sum_{l=1}^{r-1} 2^{-l-1} ( \omega_{2l-1} + \omega_{2l}) \in \left( \RR^2 \right)^{2r}.
	\end{equation}
Note that $x_i \in \RR^2$. We will argue that $(x_1, \dots, x_{2r}) = 2^{-r+1} (0_2,\ldots,0_2, \eps_{1})$.
Since $x$ is a convex combination of the elements in $\Gamma_{2,2r}$, the statement then follows from $\| \eps_1 \| = \sqrt{2}^{-1}$.

We consider $A_{2r}$ like in its construction~\eqref{eq:defA2r} as a $r \times r$ block matrix with block entries being $2 \times 2$ matrices. For $m \in [r]$ the two weights $\omega_{2m-1}$ and $\omega_{2m}$ correspond to the $m^{th}$ block row of $A_{2r}$ and have the same scalar factor in~\eqref{eq:convCombGamma}. Hence, whenever for $i \in [2r]$ the $i^{th}$ column of the $m^{th}$ block row of $A_{2k}$ contains exactly one entry equal to one (and so the other entry equals two), then the contribution of $\omega_{2m-1}$ and $\omega_{2m}$ to $x_i$ cancels due to $\eps_1 + \eps_2 = 0_2$.
In particular, in \eqref{eq:convCombGamma} all contributions of block entries equal to $B_1$ cancel. Therefore the last column of $A_{2r}$ gives
	\begin{align*}
	x_{2r} = 2^{-r} (\eps_1 + \eps_1) = 2^{-r+1} \eps_1.
	\end{align*}
Furthermore, $x_1 = x_{3} = \ldots = x_{2r-1} = 0_2$ using that also the first columns of $A_2$, of $B_2$ and of $B_3$ contain exactly one entry equal to one. For $r=1$ we are done.
If $r \geq 2$, then reading off the second column of $A_{2r}$, we find 
	\begin{align*}
	x_2 = \underbrace{2^{-2} (\eps_1 + \eps_1)}_{\text{first block row}} + \underbrace{2^{-r} (\eps_2 + \eps_2)}_{\text{last block row}} + \sum_{l=2}^{r-1} \underbrace{ 2^{-l-1} (\eps_2 + \eps_2)}_{\text{middle rows}} = 2^{-1} (\eps_1 + \eps_2) = 0_2.
	\end{align*}
Analogously, as $B_1$ does not contribute we compute for $j = 2,3,\ldots,r-1$ that
	\begin{align*}
	x_{2j} = \underbrace{2^{-j-1} (\eps_1 + \eps_1)}_{j^{th} \text{ block row}} + \underbrace{2^{-r} (\eps_2 + \eps_2)}_{\text{last block row}} + \sum_{l=j+1}^{r-1} \underbrace{ 2^{-l-1} (\eps_2 + \eps_2)}_{\text{in between rows}} = 2^{-j} (\eps_1 + \eps_2) = 0_2,
	\end{align*}
because the second columns of $B_2$ and $B_3$ are, respectively, $(2,2)^T$ and $(1,1)^T$. This proves the inequality in the case $\Gamma_{2,2r}$.

By construction, for $\Gamma_{2,2r+1}$ the same convex combination works, because the last $\RR^2$-component does not contribute as the entries of the weights alternate between $\eps_1$ and $\eps_2$.
\end{proof}

Finally, Lemma~\ref{lem:affHullQubits} and Lemma~\ref{lem:convCombQubits} together yield Theorem~\ref{thm:MarginTensor}(a), noting that for odd $d=2r+1$ one has $\; -r + 1/2 = -(d/2) + 1$.

\subsection{3-tensors}\label{sec:3tensors}

The main goal of this section is to show that the margin of $\Omega_{n,3}$ is exponentially small in $n$, i.e. to show \cref{thm:MarginTensor}(b). To do so, we set
\begin{equation}
\mathfrak{W}_n := \bigcup_{s=2}^n \lbrace (s,1,s), (s,s,1), (s-1,s,s) \rbrace \subseteq [n] \times [n] \times [n]
\end{equation}
and consider the corresponding subset
\begin{equation}
\Gamma_{n,3} := \big\lbrace (\eps_i,\eps_j,\eps_k) \mid (i,j,k) \in \mathfrak{W}_n \big \rbrace \subseteq \Omega_{n,3}.
\end{equation}

The key combinatorial idea, which is presented in the following lemma, is due to \cite[Theorem~1 with $k=0$]{krav}.\footnote{In \cite{krav} Kravtsov extensively studies so-called complete $r$-noninteger vertices ($r$-CNVs) of the three-index axial assignment polytope. For $k \in \{0,1,\ldots,n-2\}$, \cite[Theorem~1]{krav} states explicitly a $(3n-2-k)$-CNV, among these we use the $(3n-2)$-CNV (i.e. $k=0$). Moreover, \cite[Theorem~2]{krav} states that such $r$-CNVs of the three-index axial assignment polytope actually only occur for $r \in \{2n,2n+1, \ldots, 3n-2\}$, and the later theorems in \cite{krav} fully characterize the $r$-CNVs and study their combinatorial properties.}  According to \cite{krav} the special case $k=0$ is already contained in \cite[Theorem~9]{luk}.

\begin{lem}\label{lem:Kravtsov}
Let $n \geq 3$. For $(i,j,k) \in [n]^3 \setminus \big(\mathfrak{W}_n \cup \lbrace (1,1,1) \rbrace \big)$ set $\lambda_{i,j,k} := 0$. Moreover, define
\begin{align*}
\lambda_{1,1,1} := 2^{-n+1}, \quad \lambda_{1,2,2} := 1- 2^{-n+1}, \quad \lambda_{n,1,n} = \lambda_{n,n,1} := 2^{-1}
\end{align*}
and for $s=2,3,\ldots,n-1$
\begin{align*}
\lambda_{s,1,s} = \lambda_{s,s,1} := 2^{-n+s-1} ,\quad \lambda_{s,s+1,s+1} := 1-2^{-n+s} \, .
\end{align*}
Then the following equations hold:
\begin{align}
\left( \forall i \in [n]: \; \sum_{j,k = 1}^n \lambda_{i,j,k} = 1 \right), \;
\left( \forall j \in [n]: \; \sum_{i,k = 1}^n \lambda_{i,j,k} = 1 \right), \;
\left( \forall k \in [n]: \; \sum_{i,j = 1}^n \lambda_{i,j,k} = 1 \right) \, .\label{eq:marginals-krav}
\end{align}
In particular, $\,\sum_{i,j,k} \lambda_{i,j,k} = n$.
\end{lem}

\begin{proof}
This is \cite[Theorem 1 with $k=0$]{krav}. Alternatively, the statement can be checked by straightforward computation.
\end{proof}

\begin{exa}
To visualize \cref{lem:Kravtsov} it is helpful to consider the slices $\Lambda_i$ given by $(\Lambda_i)_{j,k} = \lambda_{i,j,k}$.
For $n=4$ one has
    \begin{align*}
        \Lambda_1 = {\color{blue}\frac{1}{8}} \begin{pmatrix}
            {\color{blue}1} & 0 & 0 & 0 \\ 0 & {\color{blue}7} & 0 & 0 \\ 0 & 0 & 0 & 0 \\ 0 & 0 & 0 & 0
            \end{pmatrix},\quad
        \Lambda_2 = {\color{blue}\frac{1}{8}} \begin{pmatrix}
            0 & {\color{blue}1} & 0 & 0 \\ {\color{blue}1} & 0 & 0 & 0 \\ 0 & 0 & {\color{blue}6} & 0 \\ 0 & 0 & 0 & 0
            \end{pmatrix},\quad
        \Lambda_3 = {\color{blue}\frac{1}{8}} \begin{pmatrix}
            0 & 0 & {\color{blue}2} & 0 \\ 0 & 0 & 0 & 0 \\ {\color{blue}2} & 0 & 0 & 0 \\ 0 & 0 & 0 & {\color{blue}4}
            \end{pmatrix},\quad
        \Lambda_4 = {\color{blue}\frac{1}{8}} \begin{pmatrix}
            0 & 0 & 0 & {\color{blue}4} \\ 0 & 0 & 0 & 0 \\ 0 & 0 & 0 & 0 \\ {\color{blue}4} & 0 & 0 & 0
            \end{pmatrix}.
    \end{align*}
For $n = 5$ one has
    \begin{align*}
        \Lambda_1 = {\color{blue}\frac{1}{16}} \begin{pmatrix}
            {\color{blue}1} & 0 & 0 & 0 & 0 \\ 0 & {\color{blue}15} & 0 & 0 & 0 \\ 0 & 0 & 0 & 0 & 0 \\ 0 & 0 & 0 & 0 & 0 \\ 0 & 0 & 0 & 0 & 0
            \end{pmatrix},\quad
        \Lambda_2 = {\color{blue}\frac{1}{16}} \begin{pmatrix}
            0 & {\color{blue}1} & 0 & 0 & 0 \\ {\color{blue}1} & 0 & 0 & 0 & 0 \\ 0 & 0 & {\color{blue}14} & 0 & 0 \\ 0 & 0 & 0 & 0 & 0 \\ 0 & 0 & 0 & 0 & 0
            \end{pmatrix},\quad
        \Lambda_3 = {\color{blue}\frac{1}{16}} \begin{pmatrix}
            0 & 0 & {\color{blue}2} & 0 & 0 \\ 0 & 0 & 0 & 0 & 0 \\ {\color{blue}2} & 0 & 0 & 0 & 0 \\ 0 & 0 & 0 & {\color{blue}12} & 0 \\ 0 & 0 & 0 & 0 & 0
            \end{pmatrix}\\
        \Lambda_4 = {\color{blue}\frac{1}{16}} \begin{pmatrix}
            0 & 0 & 0 & {\color{blue}4} & 0 \\ 0 & 0 & 0 & 0 & 0 \\ 0 & 0 & 0 & 0 & 0 \\ {\color{blue}4} & 0 & 0 & 0 & 0 \\ 0 & 0 & 0 & 0 & {\color{blue}8}
            \end{pmatrix},\quad
        \Lambda_5 = {\color{blue}\frac{1}{16}} \begin{pmatrix}
            0 & 0 & 0 & 0 & {\color{blue}8} \\ 0 & 0 & 0 & 0 & 0 \\ 0 & 0 & 0 & 0 & 0 \\ 0 & 0 & 0 & 0 & 0 \\ {\color{blue}8} & 0 & 0 & 0 & 0
            \end{pmatrix}.
    \end{align*}
\end{exa}


\begin{lem}\label{lem:distKravtsov}
For $n \geq 3$, it holds that $\dist\big( 0, \conv(\Gamma_{n,3}) \big) \leq 2^{-n+1}$.
\end{lem}

\begin{proof}
Define $\lambda_{i,j,k} \geq 0$ for all $i,j,k \in [n]$ as in Lemma~\ref{lem:Kravtsov}. Note that $\sum_{i = 1}^n \eps_i = 0$; thus Lemma~\ref{lem:Kravtsov} implies
	\begin{align*}
	\sum_{i,j,k} \lambda_{i,j,k} (\eps_i,\eps_j,\eps_k) = 0_{3n} \, , \quad \text{equivalently} \quad
	-2^{-n+1}(\eps_1,\eps_1,\eps_1) = \sum_{(i,j,k) \in \mathfrak{W}_n} \lambda_{i,j,k}(\eps_i,\eps_j,\eps_k).
	\end{align*}
Normalizing the latter equation we obtain
	\begin{align*}
	x:= -\frac{1}{c \, 2^{n-1}} (\eps_1,\eps_1,\eps_1) \in \conv(\Gamma_{n,3}), \quad \text{where }\;
	c := \sum_{(i,j,k)\in \mathfrak{W}_n} \lambda_{i,j,k} = n-2^{-n+1} \geq \sqrt{3}.
	\end{align*}
Finally, $\norm{\eps_1}^2 \leq 1$ implies $\norm{x} \leq c^{-1} 2^{-n+1} \sqrt{3} \leq 2^{-n+1}$.
\end{proof}

To finish the proof of Theorem~\ref{thm:MarginTensor}(b) we are left to show $0 \notin \conv(\Gamma_{n,3})$. We actually prove the stronger statement $0 \notin \aff(\Gamma_{n,3})$.

\begin{lem}\label{lem:affineHullKravtsov}
The zero vector is not contained  in the affine hull of $\Gamma_{n,3}$.
\end{lem}

\begin{proof}
For a proof by contradiction we assume $0 \in \aff(\Gamma_{n,3})$.  Then there exist $a_s, b_s, c_s \in \RR$ for $s=2,3,\ldots,n$ such that $\sum_s a_s + b_s + c_s = 1$ and
	\begin{align*}
	\sum_{s=2}^n \big( \, a_s (\eps_s, \eps_1, \eps_s) + b_s (\eps_s, \eps_s, \eps_1) + c_s (\eps_{s-1}, \eps_s, \eps_s) \, \big) = (0_n, 0_n, 0_n) \in (\RR^n)^3.
	\end{align*}
In each of the three $\RR^n$-components we obtain $0_n$ as an affine linear combination of $\eps_1, \ldots,\eps_n$.
Applying Lemma~\ref{lem:convCombEps-i} to the coefficient of $\eps_{s-1}$ in the first component,  respectively to the coefficient of $\eps_s$ in the second and third component yields
	\begin{align}
	a_{s-1} + b_{s-1} + c_{s} &= \frac{1}{n} \quad \text{ for } s=2,3,\ldots,n \label{eq:prpAffHull1} \\
	\text{respectively } \qquad
	b_s+ c_s = a_s + c_s &= \frac{1}{n} \quad \text{ for } s=2,3,\ldots,n  \qquad\qquad \label{eq:prpAffHull2}
 	\end{align}
where we necessarily set $a_1 = b_1 := 0$. Equation~\eqref{eq:prpAffHull1} for $s=2$ is $c_2 = n^{-1}$ and hence $a_2=b_2=0$ by \eqref{eq:prpAffHull2} for $s=2$. But now \eqref{eq:prpAffHull1} for $s=3$ gives $c_3 = n^{-1}$ and we can proceed inductively to conclude $c_s = n^{-1}$ and $a_s = b_s = 0$ for all $s=2,3,\ldots,n$.  This gives the contradiction $1 = \sum_{s=2}^n (a_s + b_s + c_s) = \frac{n-1}{n}$, so we must have $0 \notin \aff(\Gamma_{n,3})$. Another contradiction arises when one applies \cref{lem:convCombEps-i} to the coefficient $\eps_n$ in the first component, which yields $a_n + b_n = n^{-1}$.
\end{proof}



\subsection{$d$-tensors}\label{subsec:dTensors}
In this subsection we show that the margin of $\Omega_{n,d}$ is inverse exponential in $nd$ for $n, d \geq 3$, proving part $(c)$ of \cref{thm:MarginTensor}.

Let us give some intuition for our construction. The main idea is to recycle the construction from the previous subsection for some multiple of $n$, i.e. considering $\mathfrak{W}_{rn}$ for $r \geq 2$. Thereby, the main challenge is to ensure that the constructed subset of $\Omega_{n,d}$ does not contain zero in its convex hull. We can try to extend the elements of $\Omega_{n,3}$ to elements of $\Omega_{n,d}$. One natural idea is duplicate each component $d/3$ times, i.e. when $d=6$ the vector $(\eps_i, \eps_j, \eps_k) \in \Omega_{n,3}$ becomes $(\eps_i, \eps_i, \eps_j, \eps_j, \eps_k, \eps_k) \in \Omega_{n,6}$. However, we need a subset of $\Omega_{n,d}$ with $rn$ many elements to imitate the construction from the previous subsection. We still extend the elements of $\Omega_{n,3}$ in this way, but will additionally ``shift'' and ``twist'' by some functions $\sigma_1, \dots, \sigma_{2r-1} \colon [rn] \to [n]$, so that the elements of our set will look like 
$$
\left( \eps_{\sigma_1(i)}, \ldots, \eps_{\sigma_{d/3}(i)},
	\eps_{\sigma_1(j)}, \ldots, \eps_{\sigma_{d/3}(j)}, \eps_{\sigma_1(k)}, \ldots, \eps_{\sigma_{d/3}(k)} \right)
  $$
for $d/3 = 2r-1$ and $(i,j,k)$ in $\mathfrak{W}_{rn}$. We now set about choosing the functions $\sigma_k$. For this, let $n \geq 3$ and fix a natural number $r \geq 2$. It is convenient to use an \emph{adjusted} modulo $n$ function $\mathrm{mod}' \;\; n$ that takes values in $[n]$, i.e. instead of zero it outputs $n$.
For $i \in [r]$ we consider
	\begin{align*}
	&\sigma_{i} \colon [r n] \to [n], \quad j \mapsto \left\lceil \frac{j + (i-1)}{r} \right\rceil \quad \mathrm{mod}' \;\; n \\
	&\sigma_{r + i} := \sigma_1 \circ (r - i + 1 \quad r + 1) \colon [rn] \to [n]
	\end{align*}
where $(r - i + 1 \quad r + 1)$ denotes the corresponding transposition in the symmetric group of $[rn]$.\footnote{We stress that we always take $\sigma_1$ (and \emph{not} $\sigma_i$) to define $\sigma_{r+i}$.} We only need the first $2 r - 1$ of these functions and combine them to obtain 
	\begin{align*}
	\sigma \colon [r n] \to [n]^{2 r - 1}, \quad j \mapsto \big( \sigma_1(j), \sigma_2(j),\ldots, \sigma_{2 r - 1}(j) \big).
	\end{align*}

\begin{exa}\label{exa:sigmaCase3}
For $r = 3$ the functions $\sigma_1, \sigma_2, \ldots, \sigma_6$ are sketched by the following table.
	\begin{center}
	\begin{tabular}{|c||c|c|c|c|c|c|c|c|c|c|c|c|c|}
	\hline 
	$j$ & $1$ & $2$ & $3$ & $4$ & $5$ & $6$ & $\cdots$ & $3n -5$ & $3n - 4$ & $3n-3$ & $3n -2$ & $3n - 1$ & $3n$ \\ 
	\hline \hline
	$\sigma_1$ & \cellcolor{cyan}$1$ & \cellcolor{cyan}$1$ & \cellcolor{cyan}$1$ & \cellcolor{Yellow}$2$ & \cellcolor{Yellow}$2$ & \cellcolor{Yellow}$2$ & $\cdots$ & \cellcolor{YellowGreen}$n-1$ & \cellcolor{YellowGreen}$n-1$ & \cellcolor{YellowGreen}$n-1$ & \cellcolor{YellowOrange}$n$ & \cellcolor{YellowOrange}$n$ & \cellcolor{YellowOrange}$n$ \\
	\hline 
	$\sigma_2$ & \cellcolor{cyan}$1$ & \cellcolor{cyan}$1$ & \cellcolor{Yellow}$2$ & \cellcolor{Yellow}$2$ & \cellcolor{Yellow}$2$ & \cellcolor{Tan!70}$3$ & $\cdots$ & \cellcolor{YellowGreen}$n-1$ & \cellcolor{YellowGreen}$n-1$ & \cellcolor{YellowOrange}$n$ & \cellcolor{YellowOrange}$n$ & \cellcolor{YellowOrange}$n$ & \cellcolor{cyan}$1$ \\ 
	\hline 
	$\sigma_3$ & \cellcolor{cyan}$1$ & \cellcolor{Yellow}$2$ & \cellcolor{Yellow}$2$ & \cellcolor{Yellow}$2$ & \cellcolor{Tan!70}$3$ & \cellcolor{Tan!70}$3$ & $\cdots$ & \cellcolor{YellowGreen}$n-1$ & \cellcolor{YellowOrange}$n$ & \cellcolor{YellowOrange}$n$ & \cellcolor{YellowOrange}$n$ & \cellcolor{cyan}$1$ & \cellcolor{cyan}$1$ \\ 
	\hline \hline
	$\sigma_4$ & \cellcolor{cyan}$1$ & \cellcolor{cyan}$1$ & \cellcolor{Yellow}$2$ & \cellcolor{cyan}$1$ & \cellcolor{Yellow}$2$ & \cellcolor{Yellow}$2$ & $\cdots$ & \cellcolor{YellowGreen}$n-1$ & \cellcolor{YellowGreen}$n-1$ & \cellcolor{YellowGreen}$n-1$ & \cellcolor{YellowOrange}$n$ & \cellcolor{YellowOrange}$n$ & \cellcolor{YellowOrange}$n$ \\  
	\hline 
	$\sigma_5$ & \cellcolor{cyan}$1$ & \cellcolor{Yellow}$2$ & \cellcolor{cyan}$1$ & \cellcolor{cyan}$1$ & \cellcolor{Yellow}$2$ & \cellcolor{Yellow}$2$ & $\cdots$ & \cellcolor{YellowGreen}$n-1$ & \cellcolor{YellowGreen}$n-1$ & \cellcolor{YellowGreen}$n-1$ & \cellcolor{YellowOrange}$n$ & \cellcolor{YellowOrange}$n$ & \cellcolor{YellowOrange}$n$ \\ 
	\hline
	$\sigma_6$ & \cellcolor{Yellow}$2$ & \cellcolor{cyan}$1$ & \cellcolor{cyan}$1$ & \cellcolor{cyan}$1$ & \cellcolor{Yellow}$2$ & \cellcolor{Yellow}$2$ & $\cdots$ & \cellcolor{YellowGreen}$n-1$ & \cellcolor{YellowGreen}$n-1$ & \cellcolor{YellowGreen}$n-1$ & \cellcolor{YellowOrange}$n$ & \cellcolor{YellowOrange}$n$ & \cellcolor{YellowOrange}$n$ \\ 
	\hline 
	\end{tabular} 
	\end{center}
For $r = 3$ and $n = 5$ the functions $\sigma_1, \sigma_2, \ldots, \sigma_6$ are given by the following table.    
    \begin{center}
	\begin{tabular}{|*{16}{>{\centering\arraybackslash}p{15pt}|}}
	\hline 
	$j$ & $1$ & $2$ & $3$ & $4$ & $5$ & $6$ & $7$ & $8$ & $9$ & $10$ & $11$ & $12$ & $13$ & $14$ & $15$ \\ 
	\hline \hline
	$\sigma_1$ & \cellcolor{cyan}$1$ & \cellcolor{cyan}$1$ & \cellcolor{cyan}$1$ & \cellcolor{Yellow}$2$ & \cellcolor{Yellow}$2$ & \cellcolor{Yellow}$2$ & \cellcolor{Tan!70}$3$ & \cellcolor{Tan!70}$3$ & \cellcolor{Tan!70}$3$ & \cellcolor{YellowGreen}$4$ & \cellcolor{YellowGreen}$4$ & \cellcolor{YellowGreen}$4$ & \cellcolor{YellowOrange}$5$ & \cellcolor{YellowOrange}$5$ & \cellcolor{YellowOrange}$5$ \\
	\hline 
	$\sigma_2$ & \cellcolor{cyan}$1$ & \cellcolor{cyan}$1$ & \cellcolor{Yellow}$2$ & \cellcolor{Yellow}$2$ & \cellcolor{Yellow}$2$ & \cellcolor{Tan!70}$3$ & \cellcolor{Tan!70}$3$ & \cellcolor{Tan!70}$3$ & \cellcolor{YellowGreen}$4$ & \cellcolor{YellowGreen}$4$ & \cellcolor{YellowGreen}$4$ & \cellcolor{YellowOrange}$5$ & \cellcolor{YellowOrange}$5$ & \cellcolor{YellowOrange}$5$ & \cellcolor{cyan}$1$ \\ 
	\hline 
	$\sigma_3$ & \cellcolor{cyan}$1$ & \cellcolor{Yellow}$2$ & \cellcolor{Yellow}$2$ & \cellcolor{Yellow}$2$ & \cellcolor{Tan!70}$3$ & \cellcolor{Tan!70}$3$ & \cellcolor{Tan!70}$3$ & \cellcolor{YellowGreen}$4$ & \cellcolor{YellowGreen}$4$ & \cellcolor{YellowGreen}$4$ & \cellcolor{YellowOrange}$5$ & \cellcolor{YellowOrange}$5$ & \cellcolor{YellowOrange}$5$ & \cellcolor{cyan}$1$ & \cellcolor{cyan}$1$ \\ 
	\hline \hline
	$\sigma_4$ & \cellcolor{cyan}$1$ & \cellcolor{cyan}$1$ & \cellcolor{Yellow}$2$ & \cellcolor{cyan}$1$ & \cellcolor{Yellow}$2$ & \cellcolor{Yellow}$2$ & \cellcolor{Tan!70}$3$ & \cellcolor{Tan!70}$3$ & \cellcolor{Tan!70}$3$ & \cellcolor{YellowGreen}$4$ & \cellcolor{YellowGreen}$4$ & \cellcolor{YellowGreen}$4$ & \cellcolor{YellowOrange}$5$ & \cellcolor{YellowOrange}$5$ & \cellcolor{YellowOrange}$5$ \\
	\hline 
	$\sigma_5$ & \cellcolor{cyan}$1$ & \cellcolor{Yellow}$2$ & \cellcolor{cyan}$1$ & \cellcolor{cyan}$1$ & \cellcolor{Yellow}$2$ & \cellcolor{Yellow}$2$ & \cellcolor{Tan!70}$3$ & \cellcolor{Tan!70}$3$ & \cellcolor{Tan!70}$3$ & \cellcolor{YellowGreen}$4$ & \cellcolor{YellowGreen}$4$ & \cellcolor{YellowGreen}$4$ & \cellcolor{YellowOrange}$5$ & \cellcolor{YellowOrange}$5$ & \cellcolor{YellowOrange}$5$ \\
	\hline
	$\sigma_6$ & \cellcolor{Yellow}$2$ & \cellcolor{cyan}$1$ & \cellcolor{cyan}$1$ & \cellcolor{cyan}$1$ & \cellcolor{Yellow}$2$ & \cellcolor{Yellow}$2$ & \cellcolor{Tan!70}$3$ & \cellcolor{Tan!70}$3$ & \cellcolor{Tan!70}$3$ & \cellcolor{YellowGreen}$4$ & \cellcolor{YellowGreen}$4$ & \cellcolor{YellowGreen}$4$ & \cellcolor{YellowOrange}$5$ & \cellcolor{YellowOrange}$5$ & \cellcolor{YellowOrange}$5$ \\
	\hline 
	\end{tabular} 
	\end{center}
\end{exa}

\begin{rem}\label{rem:Sigma}
By construction, each element of $[n]$ is attained exactly $r$-times by $\sigma_k$, $k \in [2 r- 1]$. Moreover, the definition of $\sigma_1, \ldots, \sigma_r$ yields that $\sigma$ is injective. \end{rem}

For $i,j,k \in [r n]$ we introduce the short-hand
	\begin{align*}
	\eps_{\sigma(i)} &:= \left( \eps_{\sigma_1(i)}, \eps_{\sigma_{2}(i)}, \ldots, \eps_{\sigma_{2 r - 1}(i)} \right) \in \left( \RR^n \right)^{2 r - 1} \\
	\eps_{\sigma(i), \sigma(j), \sigma(k)} &:= \left( \eps_{\sigma_1(i)}, \ldots, \eps_{\sigma_{2 r - 1}(i)},
	\eps_{\sigma_1(j)}, \ldots, \eps_{\sigma_{2 r - 1}(j)}, \eps_{\sigma_1(k)}, \ldots, \eps_{\sigma_{2 r - 1}(k)} \right)
	\in \left( \RR^n \right)^{6r - 3}
	\end{align*}
and we set\footnote{One could suggest to consider the set $\lbrace \eps_{\sigma(i), \sigma(j), \sigma(k)} \mid (i,j,k) \in \mathfrak{W}_{r n} \rbrace$, but this still won't ensure that zero is not in the convex hull. The intuition behind is, that $\Gamma_{n,3}$ from the last section is ``nearly at the limit'', i.e. $0 \notin \conv(\Gamma_{n,3})$ but $0 \in \conv(\Gamma_{n,3} \cup \{ (\eps_1,\eps_1,\eps_1) \})$. Now the function $\sigma$ ``introduces $2r-2$ additional linear relations'' as $\eps_{\sigma(i)} \in (\id_n^\perp)^{2r-1}$, since the orthogonal complement $\id_n^\perp \subseteq \RR^n$ has codimension one while $(\id_n^\perp)^{2r-1} \subseteq (\RR^n)^{2r-1}$ has codimension $2r-1$. Thus, it is reasonable to remove $2r-2$ many elements from $\mathfrak{W}_{rn}$.}
	\begin{align*}
	\mathfrak{J}_r := \big\lbrace (s,1,s), (s,s,1) \mid s = 2,3,\ldots, r \big\rbrace \subseteq \ZZ^3.
	\end{align*}
In the following we show that the convex hull of the set
	\begin{align*}
	\Gamma_{n, 6 r - 3} = \big\lbrace \eps_{\sigma(i), \sigma(j), \sigma(k)} \mid (i,j,k) \in \mathfrak{W}_{r n} \setminus \mathfrak{J}_{r} \big\rbrace \subseteq \Omega_{n, 6 r - 3} \subseteq \Big( \big( \RR^n \big)^{2 r - 1} \Big)^3
	\end{align*}
does not contain the zero vector, but is very close to it.

\begin{lem}\label{lem:convStackingKravtsov}
For $n \geq 3$ and $r \geq 2$ it holds that $0 \notin \aff \left( \Gamma_{n, 6 r - 3} \right)$.
\end{lem}

Below we give the proof in the special case $r=3$, in which all main ideas of the general proof become apparent and visible. The proof for the general statement is given in \cref{sec:StackingKravtsovGeneral} and certainly looks technical at a first encounter. Therefore, we strongly suggest that the reader first reads the proof for $r=3$ below. 



\begin{proof}[Proof of \cref{lem:convStackingKravtsov} for $r=3$]
For the sake of contradiction assume that $0 \in \aff(\Gamma_{n, 15})$. Then there are coefficients $a_s, b_s, c_s \in \RR$, where $2 \leq s \leq 3 n$, such that $a_2 = a_3 = b_2 = b_3 = 0$, $\sum_s (a_s + b_s + c_s) = 1$ and
	\begin{align}
	\sum_{s= 2}^{3 n} \left( a_s \, \eps_{\sigma(s),\sigma(1),\sigma(s)}
	+ b_s \, \eps_{\sigma(s),\sigma(s),\sigma(1)} + c_s \, \eps_{\sigma(s-1),\sigma(s),\sigma(s)}  \right) = 0 \in (\RR^n)^{15}.\label{eq:main-comboCase3}
	\end{align}
	The bulk of our work will consist of proving the equations
	\begin{align}\label{eq:StackBandCCase3}
	b_2 + c_2 &= b_3 + c_3 = \ldots = b_{3 n} + c_{3 n}\\
	\label{eq:StackAandCCase3} a_2 + c_2 &= a_3 + c_3 = \ldots = a_{3 n} + c_{3 n}.
	\end{align}
	From here we will derive a contradiction. We now set about proving \cref{eq:StackAandCCase3,eq:StackBandCCase3}. Rewrite the left-hand-side of \cref{eq:main-comboCase3} as the collection for $k \in [5]$ of the following affine linear combinations of $\eps_1,\ldots,\eps_n$ in $\RR^n$:
	\begin{align}
	\sum_{s= 2}^{3n} \left( a_s \, \eps_{\sigma_k(s)} 
	+ b_s \, \eps_{\sigma_k(s)} + c_s \, \eps_{\sigma_k(s-1)}  \right) &= 0 \label{eq:StackComp1Case3}\\
	\sum_{s= 2}^{3n} \left( a_s \, \eps_{\sigma_k(1)}
	+ b_s \, \eps_{\sigma_k(s)} + c_s \, \eps_{\sigma_k(s)}  \right) &= 0 \label{eq:StackComp2Case3} \\
	\sum_{s= 2}^{3n} \left( a_s \, \eps_{\sigma_k(s)}
	+ b_s \, \eps_{\sigma_k(1)} + c_s \, \eps_{\sigma_k(s)}  \right) &= 0. \label{eq:StackComp3Case3}
	\end{align}
If we expand each expression as an affine linear combination of the $\eps_l$, then by Lemma~\ref{lem:convCombEps-i} the coefficient of $\eps_l$ must be $n^{-1}$ for all $l \in [n]$. Translating this for equation \eqref{eq:StackComp1Case3} with $k = 2$, $l=2,\ldots,n$ and using \cref{exa:sigmaCase3} we obtain 


 \begin{align}
	(a_{m-3} + a_{m-2} + a_{m-1}) + (b_{m-3} + b_{m-2} + b_{m-1})  + (c_{m-2} + c_{m-1} + c_{m}) &= \frac{1}{n} \label{eq:First2}
		\end{align}
for $m = 6,9,12,\dots, 3n$. A similar calculation for $k=1,3$ and $l=2,\ldots,n$ shows \cref{eq:First2} holds for all $5 \leq m \leq 3n +1$, where we set  $c_{3 n + 1} := 0$.


Similarly for \cref{eq:StackComp2Case3} with $l=2,\ldots,n$ and $k=1,2,3$ we obtain for $4 \leq m \leq 3n$ that
	\begin{align}
	(b_{m-2} + c_{m-2})  + (b_{m-1} + c_{m-1}) + (b_{m} + c_m)  &= \frac{1}{n} \label{eq:Second1} 
	\end{align}
and the same equations with ``$b$'' replaced by ``$a$'' when considering \cref{eq:StackComp3Case3}.

In the following we prove \cref{eq:StackBandCCase3}. Subtracting \eqref{eq:Second1} from \eqref{eq:Second1} with values of $m$ differing by one, we deduce that 
\begin{align*}
	b_{2} + c_2 &= b_5 + c_5 = \ldots = b_{3n-1} + c_{3n-1}\\
	b_3 + c_3 &= b_6 + c_6 = \ldots = b_{3n} + c_{3n}, \\
	\qquad \text{and} \qquad b_4 + c_4 &= b_7 + c_7 = \ldots = b_{3n-2} + c_{3n-2}.
	\end{align*}
Next we deduce \cref{eq:StackBandCCase3} by showing $b_2 + c_2 = b_3 + c_3 = b_4 + c_4$. 

To do so, we apply  \cref{lem:convCombEps-i} to \eqref{eq:StackComp2Case3} for the coefficient of $\eps_2$ using \cref{exa:sigmaCase3}, which yields for $k=4,5$ the equations
	\begin{align}
	(b_3 + c_3) + (b_5 + c_5) + (b_6 + c_6)  &= \frac{1}{n}  \label{eq:finishBandC1}\\
	(b_2 + c_2) + (b_5 + c_5) + (b_6  + c_6)  &= \frac{1}{n} \label{eq:finishBandC2}
	\end{align}
respectively. Subtracting the two shows $b_2 + c_2  = b_3 + c_3$, and we have $b_3 + c_3  = b_4 + c_4$ via subtracting \eqref{eq:finishBandC1} from \eqref{eq:Second1} for $m=6$. This completes the proof of \cref{eq:StackBandCCase3}; using \cref{eq:StackComp3Case3} we similarly deduce \cref{eq:StackAandCCase3}.

To get a contradiction we show that $a_s = b_s = c_s = 0$ for all $s = 2,3,\ldots, 3 n$. For this, we set $a := \sum_{s} a_s$ and $b := \sum_s b_s$, and recall that we have defined $a_2 = a_3 = b_2 = b_3 = 0$. This time we use \cref{lem:convCombEps-i} applied to the coefficient of $\eps_1$ in \eqref{eq:StackComp1Case3}, in \eqref{eq:StackComp2Case3} and in \eqref{eq:StackComp3Case3} respectively for $k=1$ to get
	\begin{equation}\label{eq:Case3contradiction}
	c_2 + c_3 + c_4 = \frac{1}{n}, \qquad
	a + c_2 + c_3 = \frac{1}{n}  \qquad \text{ and } \qquad
	b + c_2 + c_3 = \frac{1}{n}
	\end{equation}
respectively. We deduce from these three equations that $a=b=c_4$. Furthermore, $b_2 = b_3 = 0$ shows that \eqref{eq:Second1} for $m=4$ is $b_4 + (c_2 + c_3 + c_4) = n^{-1}$. Subtracting from the latter the left-hand equation in \eqref{eq:Case3contradiction} yields $b_4 = 0$. Similarly, $a_4 = 0$ follows from $a_2 = a_3 = 0$ and the analogous equation of \eqref{eq:Second1} with $a$'s replaced by $b$'s.

Now, \eqref{eq:First2} for $m=5$ simplifies to $c_3 + c_4 + c_5 = n^{-1}$. Thus, $c_2 = c_5$ with \eqref{eq:Case3contradiction} and therefore $a_5 = b_5 = 0$ by \eqref{eq:StackBandCCase3}, \eqref{eq:StackAandCCase3} and $a_2=b_2=0$. This simplifies \eqref{eq:First2} for $m=6$ to $c_4 + c_5 + c_6 = n^{-1}$. Hence, $c_3 = c_6$ as we also have $c_3 + c_4 + c_5 = n^{-1}$ and we get via \eqref{eq:StackBandCCase3} and \eqref{eq:StackAandCCase3} that $a_6 = b_6 = 0$. The latter in turn shows that \eqref{eq:First2} for $m=7$ becomes $c_5 + c_6 + c_7 = n^{-1}$, so $c_4 = c_7$ and $a_7 = b_7 = 0$ by, again, \eqref{eq:StackBandCCase3} and \eqref{eq:StackAandCCase3}.

It should have become apparent that we can proceed inductively in the same manner with \eqref{eq:First2} for $m=5,\ldots,3n+1$; thereby using \eqref{eq:StackBandCCase3} and \eqref{eq:StackAandCCase3} to deduce $a_s = b_s = 0$ for all $s=2,3,\ldots,3n$. In particular, $a = b = c_4 = 0$. Finally, \cref{eq:StackBandCCase3} implies $c_4 = c_s$ for all $s=2,3,\ldots,3n$, which gives the desired contradiction.
\end{proof}

We finish the proof of part $(c)$ of \cref{thm:MarginTensor} by showing the following Lemma.

\begin{lem}\label{lem:distStackingKravtsov}
Let $n \geq 3$ and $r \geq 2$. Then 
	\begin{align*}
	\dist \big(0, \conv(\Gamma_{n,6 r - 3}) \big) \leq \frac{\sqrt{6}}{(n-1)\sqrt{r}} \; 2^{-r(n-1) + 1} \leq 2^{- r (n-1) + 1} .
	\end{align*}
\end{lem}

\begin{proof}
We set $N := r n$ and for $i,j,k \in [N]$ we set $\lambda_{i,j,k}$ as in \cref{lem:Kravtsov} applied for the dimension $N$. Then \cref{eq:marginals-krav} of \cref{lem:Kravtsov} yields
	\begin{align*}
	&\sum_{i,j,k=1}^N \lambda_{i,j,k} \, \left( \eps_{\sigma(i)}, \eps_{\sigma(j)}, \eps_{\sigma(k)} \right)\\
	= &\sum_{i,j,k=1}^N \lambda_{i,j,k} \, \left( \eps_{\sigma(i)}, 0, 0 \right)+ \sum_{i,j,k=1}^N \lambda_{i,j,k} \, \left( 0, \eps_{ \sigma(j)}, 0 \right) + \sum_{i,j,k=1}^N \lambda_{i,j,k} \, \left( 0, 0, \eps_{\sigma(k)} \right) \\
	= &\sum_{i=1}^N \left( \eps_{\sigma(i)}, 0, 0 \right)+ \sum_{j=1}^N  \left( 0, \eps_{ \sigma(j)}, 0 \right) + \sum_{k=1}^N \left( 0, 0, \eps_{\sigma(k)} \right) 
	= \sum_{i=1}^N \eps_{\sigma(i),\sigma(i),\sigma(i)} = 0 \in \left(\RR^n \right)^{6 r - 3},
	\end{align*}
where we used in the last step equation~\eqref{eq:SumEps-i} and \cref{rem:Sigma}, i.e. that each element of $[n]$ is attained exactly $r$-many times by all $\sigma_k \colon [r n] \to [n]$, $k \in [2 r - 1]$. Because $\mathfrak{W}_{N}$ contains the support of $\lambda$ apart from the element $(1,1,1)$, we have
	\begin{align}
	\sum_{(i,j,k) \in \mathfrak{W}_{N} \setminus \mathfrak{J}_{r}} \lambda_{i,j,k} \, \eps_{\sigma(i), \sigma(j), \sigma(k)} =  - \lambda_{1,1,1} \, \eps_{\sigma(1), \sigma(1), \sigma(1)} -  \sum_{(i,j,k) \in \mathfrak{J}_{r}} \lambda_{i,j,k} \, \eps_{\sigma(i), \sigma(j), \sigma(k)} =: x \in \left( \RR^n \right)^{6 r - 3},\label{eq:x-stacked}
	\end{align}
which is an element in the positive cone of $\Gamma_{n,6 r - 3} = \lbrace \eps_{\sigma(i), \sigma(j), \sigma(k)} \mid (i,j,k) \in \mathfrak{W}_N \setminus \mathfrak{J}_{r} \rbrace$.
Normalizing the latter equation with
	\begin{align*}
	c := \sum_{(i,j,k) \in \mathfrak{W}_{N} \setminus \mathfrak{J}_{r}} \lambda_{i,j,k} 
	= \sum_{i,j,k=1}^N \lambda_{i,j,k} - \left( \lambda_{1,1,1} + \sum_{(i,j,k)\in \mathfrak{J}_{r}} \lambda_{i,j,k} \right) \geq N-1
	\end{align*}
shows $c^{-1} x \in \conv(\Gamma_{n,6 r - 3})$. To bound the norm of $c^{-1} x$ we compute
	\begin{align*}
	\lambda_{1,1,1} + \sum_{(i,j,k) \in \mathfrak{J}_{r}} \lambda_{i,j,k} &= 2^{-N+1} + \sum_{s=2}^r \left( \lambda_{s,1,s} + \lambda_{s,s,1} \right) \\
	&= 2^{-N+1} + \sum_{s=2}^r \left( 2^{-N+s-1} + 2^{-N+s-1} \right)
	= \sum_{s=1}^r 2^{-N+s} < 2^{-N + r + 1}.
	\end{align*}
Finally, using $\| \eps_{i_1, i_2, \ldots, i_{6 r  -3}} \| \leq \sqrt{6 r - 3}$ for any $i_1, i_2, \ldots,i_{6r - 3} \in [n]$ together with the triangle inequality on \cref{eq:x-stacked} implies
	\begin{align*}
	\| c^{-1} x\| \leq \frac{\sqrt{6 r - 3}}{N - 1} \; 2^{-N + r + 1}
	\leq \frac{\sqrt{6}}{(n-1)\sqrt{r}} \; 2^{-N + r + 1} \leq 2^{-N + r + 1} = 2^{- r (n-1) + 1},
	\end{align*}
where we used $n \geq 3$ and $r \geq 2$ for $\sqrt{6} \leq (n-1) \sqrt{r}$.
\end{proof}







































































































































































































































































































































































































































































































































































































































































































































































































































































































































































































































































































































































\appendix\setcounter{section}{2}\section{Original Proposition C.1, part (1): padding the commutative margin}\label{subsec:extension}	\begin{align*}
	\pi_{n,d} \colon \SL(n)^d \to \GL\left( (\CC^n)^{\otimes d} \right), \; (g_1, \ldots, g_d) \mapsto g_1 \otimes \cdots \otimes g_d,
	\end{align*}
which set of weights is $\Omega(\pi_{n,d}) = \Omega_{n,d} = \lbrace \eps_{i} \mid i \in [n] \rbrace^d \subseteq (\RR^n)^d$.

\begin{prp}[Original Proposition C.1, part (1)]\label{prp:dTensorsPadding}
Let $G:= \SL(n)^d$ and $n,d \geq 1$. Consider a set of weights $\Gamma_{n,d} \subseteq \Omega_{n,d}$ such that $0 \notin \conv(\Gamma_{n,d})$, i.e. $\Gamma_{n,d}$ witnesses the inequality $\gamma(\Omega_{n,d}) = \gamma_{\T}(\pi_{n,d}) \leq \dist (0, \conv(\Gamma_{n,d}))$.
	\begin{enumerate}
	\item  Then $\gamma(\Omega_{n,d+1}) \leq \dist \big(0, \conv(\Gamma_{n,d}) \big)$. Consequently, $\gamma(\Omega_{n,d+1}) \leq \gamma(\Omega_{n,d})$.

	\end{enumerate}

\end{prp}

\begin{proof}
To prove the statement we set for $r \geq 1 $
	\[ \Delta_r := \{ (\eps_i,\ldots,\eps_i) \mid i \in [n] \} \subseteq (\RR^n)^r \qquad \text{and} \qquad
	\Gamma_{n,d+r} := \Gamma_{n,d} \times \Delta_r \subseteq \Omega(\pi_{n,d+r}) . \]
By \cref{eq:SumEps-i} we have $0 \in \conv(\Delta_r)$ and therefore
	\begin{align*}
	\conv(\Gamma_{n,d+r}) = \conv(\Gamma_{n,d}) \times \conv(\Delta_r)
	\supseteq \conv(\Gamma_{n,d}) \times \{0\}.
	\end{align*}
The latter implies
	\begin{equation}\label{eq:Padding}
	\dist \big( 0, \conv(\Gamma_{n,d+r}) \big) \leq \dist \big( 0, \conv(\Gamma_{n,d}) \big).
	\end{equation}
Clearly, $0 \in \conv(\Gamma_{n,d+r})$ implies $0 \in \conv(\Gamma_{n,d})$ or, by contraposition, the assumption $0 \notin \conv(\Gamma_{n,d})$ yields $0 \notin \conv(\Gamma_{n,d+r})$. The latter for $r=1$ shows $\gamma_{\T}(\pi_{n,d+1}) \leq \dist \big(0, \conv(\Gamma_{n,d + 1}) \big)$ and we conclude the first assertion with \cref{eq:Padding}.
\end{proof}

























\setcounter{section}{3}\section{Proof of \cref{lem:convStackingKravtsov}}\label{sec:StackingKravtsovGeneral}

\begin{proof}
For the sake of contradiction assume that $0 \in \aff(\Gamma_{n, 6 r - 3})$. Then there are coefficients $a_s, b_s, c_s \in \RR$, where $2 \leq s \leq r n$, such that $a_2 = \ldots = a_r = b_2 = \ldots = b_r = 0$, $\sum_s (a_s + b_s + c_s) = 1$ and
	\begin{align}
	\sum_{s= 2}^{r n} \left( a_s \, \eps_{\sigma(s),\sigma(1),\sigma(s)}
	+ b_s \, \eps_{\sigma(s),\sigma(s),\sigma(1)} + c_s \, \eps_{\sigma(s-1),\sigma(s),\sigma(s)}  \right) = 0 \in (\RR^n)^{6 r - 3}.\label{eq:main-combo}
	\end{align}
	The bulk of our work will consist of proving the equations
	\begin{align}\label{eq:StackBandC}
	b_2 + c_2 &= b_3 + c_3 = \ldots = b_{r n} + c_{r n}\\
	\label{eq:StackAandC} a_2 + c_2 &= a_3 + c_3 = \ldots = a_{r n} + c_{r n}.
	\end{align}
	From here we will derive a contradiction. We now set about proving \cref{eq:StackAandC,eq:StackBandC}. Rewrite the left-hand-side of \cref{eq:main-combo} as the collection for $k \in [2 r - 1]$ of the following affine linear combinations of $\eps_1,\ldots,\eps_n$ in $\RR^n$:
	\begin{align}
	\sum_{s= 2}^{r n} \left( a_s \, \eps_{\sigma_k(s)} 
	+ b_s \, \eps_{\sigma_k(s)} + c_s \, \eps_{\sigma_k(s-1)}  \right) &= 0 \label{eq:StackComp1}\\
	\sum_{s= 2}^{r n} \left( a_s \, \eps_{\sigma_k(1)}
	+ b_s \, \eps_{\sigma_k(s)} + c_s \, \eps_{\sigma_k(s)}  \right) &= 0 \label{eq:StackComp2} \\
	\sum_{s= 2}^{r n} \left( a_s \, \eps_{\sigma_k(s)}
	+ b_s \, \eps_{\sigma_k(1)} + c_s \, \eps_{\sigma_k(s)}  \right) &= 0. \label{eq:StackComp3}
	\end{align}
If we expand this expressions as affine linear combinations of the $\eps_l$, then by Lemma~\ref{lem:convCombEps-i} the coefficient of $\eps_l$ must be $n^{-1}$ for all $l \in [n]$. Translating this for equations \eqref{eq:StackComp1}, \eqref{eq:StackComp2} and \eqref{eq:StackComp3} respectively with $2 \leq l \leq n$ and $k \in [r]$, and using for $j \in [r]$ that
	\begin{align}
	\sigma_{k} \big(r(l-1)+j - k + 1 \big) = \left\lceil \frac{(r(l-1)+j - k + 1) + (k-1)}{r} \right\rceil = l \label{eq:sigma-inverse}
	\end{align}
we get
	\begin{align}
	&\forall \, k \in [r], l \in \{2,3,\ldots,n\} \colon  &\sum_{j=1}^r \big( a_{r(l-1)+j - k + 1} + b_{r(l-1)+j - k + 1} + c_{r(l-1)+j- k + 2} \big) &= \frac{1}{n} \label{eq:Coeff1} \\
	&\forall \, k \in [r], l \in \{2,3,\ldots,n\} \colon &\sum_{j=1}^{r} \big( b_{r(l-1)+j - k + 1} + c_{r(l-1)+j - k + 1} \big) &= \frac{1}{n} \label{eq:Coeff2} \\
	&\forall \, k \in [r], l \in \{2,3,\ldots,n\} \colon &\sum_{j=1}^{r} \big( a_{r(l-1)+j - k + 1} + c_{r(l-1)+j - k + 1} \big) &= \frac{1}{n} \label{eq:Coeff3}
	\end{align}
respectively, where we set  $c_{r n + 1} := 0$. Fixing some $l \geq 2$ and subtracting \cref{eq:Coeff2} with $k = 1$ from \cref{eq:Coeff2} for $k= 2$, we find a telescoping sum that reduces to $b_{r(l-1)} + c_{r(l-1)} = b_{r l} + c_{r l}$. Indeed, subtracting the two yields
\begin{align*}
	0 &= \sum_{j=1}^{r} \big( b_{r(l-1)+j - 1} + c_{r(l-1)+j - 1} \big) - \sum_{j=1}^{r} \big( b_{r(l-1)+j} + c_{r(l-1)+j} \big) \\
	&= \sum_{j=0}^{r-1} \big( b_{r(l-1)+j } + c_{r(l-1)+j } \big)  - \sum_{j=1}^{r} \big( b_{r(l-1)+j} + c_{r(l-1)+j} \big) \\
	&= ( b_{r(l-1)} + c_{r(l-1)}) - (b_{r l} + c_{r l}).
	\end{align*}
More generally, for $k \in  [r-1]$ combining \eqref{eq:Coeff2} for $k$ and $k \leftarrow k + 1$, implies 
	$b_{r l - k + 1} + c_{r l - k + 1} = b_{r(l-1)- k + 1} + c_{r(l-1)- k + 1}$
for all $l = 2,\ldots,n$, i.e. for every $k \in [r - 1]$ we have
	\begin{equation}\label{eq:StackBandC1}
	c_{r - k + 1} = b_{r - k + 1} + c_{r - k+ 1} = b_{2r - k + 1} + c_{2r - k + 1} = \ldots = b_{r n - k + 1} + c_{r n - k + 1}.
	\end{equation}
We are still missing the value $k = 0$, or the equations 
\begin{equation}\label{eq:StackBandC2}
	b_{r+1} + c_{r+1} = b_{2r + 1} + c_{2r + 1} = \ldots = b_{r(n-1) + 1} + c_{r(n-1) + 1}.
	\end{equation}
We obtain this by subtracting, for $l = 2, \dots, n$,  \eqref{eq:Coeff2} for $k = 1$ and $l$ from \eqref{eq:Coeff2} with $k = r$ and $l \leftarrow l + 1$ . Indeed, 
\begin{align*}0 &= \sum_{j=1}^{r} \big( b_{r l+j - r + 1} + c_{r l +j - r + 1} \big) - \sum_{j=1}^{r} \big( b_{r(l-1)+j} + c_{r(l-1)+j} \big) \\
&= \sum_{j=2}^{r+1} \big( b_{r(l-1)+j } + c_{r(l-1) +j}\big) - \sum_{j=1}^{r} \big( b_{r(l-1)+j} + c_{r(l-1)+j} \big) \\
&= \big( b_{r l + 1} + c_{r l + 1}\big) - \big( b_{r(l-1)+1} + c_{r(l-1)+1} \big).
\end{align*}

	
Lastly, we are missing the equations $b_2 + c_2 = b_3 + c_3 = \ldots = b_{r+1} + c_{r+1}$  for \cref{eq:StackBandC}. We have not yet used in \cref{eq:StackComp2} the values $k = r + m$ with $m \in [r-1]$. For this we note that
\begin{align*}
	\sigma_{r + m} \big(j \big) =  2 \quad &\text{ for } \; j \in  \{r - m + 1\} \cup \{r + 2, r + 3, \dots, 2 r\}.
	\end{align*}
We use this equation to apply Lemma~\ref{lem:convCombEps-i} to \eqref{eq:StackComp2} for $\eps_2$ and $k=r + m$ with $m \in [r-1]$ to obtain 
$$ b_{r - m + 1} + c_{r - m +1} + \sum_{j=2}^{r} \big( b_{r + j} + c_{r + j} \big) = \frac{1}{n}.
$$
We need one more equation to eliminate the right-hand term, so we use the following. Lemma~\ref{lem:convCombEps-i} applied to 
equation~\eqref{eq:Coeff2} for $k=1$ and $l=2$ yields
	\begin{align*}
	\sum_{j=1}^r \big( b_{r + j} + c_{r + j} \big) = \frac{1}{n}.
	\end{align*}
Subtracting this equation from the previous one yields, $b_{r - m + 1} + c_{r - m +1} = b_{r+1} + c_{r + 1}$ for all $m = 1,\ldots,r-1$. Together with the equations \eqref{eq:StackBandC1} and \eqref{eq:StackBandC2} we conclude \cref{eq:StackBandC}. Analogously, \eqref{eq:StackComp3} and \eqref{eq:Coeff3} can be used to obtain
	\cref{eq:StackAandC}. 

To get a contradiction we show that $a_s = b_s = c_s = 0$ for all $s = 2,3,\ldots, r n$. For this, we set $a := \sum_{s} a_s$ and $b := \sum_s b_s$. \cref{eq:sigma-inverse} still applies for $l =1, k = 1$, so Lemma~\ref{lem:convCombEps-i} applied to the coefficient of $\eps_1$ in \eqref{eq:StackComp1}, in \eqref{eq:StackComp2} and in \eqref{eq:StackComp3} respectively for $k=1$ gives
	\begin{align*}
	\sum_{j=1}^{r} c_{j+1} = \frac{1}{n}, \qquad a + \sum_{j=1}^{r-1} c_{j+1} = \frac{1}{n} \qquad \text{ and } \qquad
	b + \sum_{j=1}^{r-1} c_{j+1} = \frac{1}{n}
	\end{align*}
respectively. Subtracting the second equation from the first gives $a=c_{r+1}$, and reasoning analogously for the third yields $a=b=c_{r+1}$. Moreover, \eqref{eq:Coeff2} with $k=r$ and $l=2$ is $\sum_{j=1}^r (b_{j+1} + c_{j+1}) = n^{-1}$. Using the latter together with $b_2 = \ldots = b_r = 0$ and $\sum_{j=1}^{r} c_{j+1} = n^{-1}$ yields $b_{r+1} = 0$ and similarly $a_{r+1} = 0$ via \eqref{eq:Coeff3} with $k=r$ and $l=2$.

Since now also $a_{r + 1} = b_{r +1} = 0$, the equation~\eqref{eq:Coeff1} with $k=r$ and $l=2$ simplifies to $\sum_{j=1}^r c_{j+2} = n^{-1}$. In conjunction with $\sum_{j=1}^{r} c_{j+1} = n^{-1}$ we deduce $c_2 = c_{r+2}$ and hence $b_{r+2} = 0 = a_{r +2} $ by \eqref{eq:StackBandC} and \eqref{eq:StackAandC}. But now \eqref{eq:Coeff1} with $k=r-1$ and $l=2$ is $\sum_{j=1}^r c_{j+3} = n^{-1}$ and together with $\sum_{j=1}^r c_{j+2} = n^{-1}$ we get $c_3 = c_{r+3}$. Continuing inductively we obtain 
	\begin{align*}
	\forall \, j \in [r] \colon \quad c_{j+1} = c_{r + j + 1} \quad \text{ and } \quad a_{r + j +1} = b_{r + j + 1} = 0
	\end{align*}
via \eqref{eq:Coeff1} with $l=2$, $k \in [r]$ and via \eqref{eq:StackBandC}, \eqref{eq:StackAandC}. Then \eqref{eq:Coeff1} with $k=r$ and $l=3$ simplifies to $\sum_{j=1}^r c_{r + j+2} = n^{-1}$ and together with $n^{-1} = \sum_{j=1}^{r} c_{j+1} =  \sum_{j=1}^{r} c_{r + j +1}$ we have $c_{r + 2} = c_{2r + 2}$. Hence, $b_{2 r+2} = 0 = a_{2 r + 2}$ via \eqref{eq:StackBandC} respectively \eqref{eq:StackAandC}. Continuing inductively in the outlined manner with equation~\eqref{eq:Coeff1} for $k \in [r]$, $l=3,\ldots,n$ and with the equations \eqref{eq:StackBandC} and \eqref{eq:StackAandC} we conclude $a_s = b_s = 0$ for all $s=2,3 \ldots, r n$, so $a=b=0$. Finally, \eqref{eq:StackBandC} implies $c_{r + 1} = c_s$ for all $s = 2,\ldots, r n$, but $c_{r + 1} = b = 0$ giving the desired contradiction.
\end{proof}




























































